\documentclass[10pt,a4paper]{amsart}
\usepackage[margin=19mm]{geometry}
\usepackage{amsmath,amssymb,mathtools}
\usepackage[scr=rsfs,cal=euler]{mathalfa}
\usepackage{xfrac}
\usepackage[unicode,hidelinks,bookmarks=false]{hyperref}
\newcommand{\ie}{i.e.\ }
\newcommand{\cf}{cf.\ }
\newcommand{\resp}{respectively\ }
\newcommand{\Subsection}{\subsection}
\providecommand{\nobreakdash}{\nobreak\hbox{-}}
\setlength{\emergencystretch}{2em}
\multlinegap0pt
\usepackage{bm}
\usepackage{bbm}
\usepackage{dsfont}
\usepackage{xspace}
\usepackage{cancel}
\usepackage{mathtools}
\usepackage{cleveref}
\usepackage{amsthm}
\makeatletter
\@ifundefined{theorem}{\newtheorem{theorem}{Theorem}}{}
\@ifundefined{fact}{\newtheorem{fact}[theorem]{Fact}}{}
\@ifundefined{lemma}{\newtheorem{lemma}[theorem]{Lemma}}{}
\makeatother
\newcommand{\EE}{\mathbb{E}}
\newcommand{\E}{\mathbb{E}}
\newcommand{\PP}{\mathbb{P}}
\newcommand{\cB}{\mathcal{B}}
\newcommand{\cC}{\mathcal{C}}
\newcommand{\cD}{\mathcal{D}}
\newcommand{\cZ}{\mathcal{Z}}
\title[The largest $K\_r$-free set of vertices in a random graph]{The largest $K\_r$-free set of vertices in a random graph}
\author{Tom Bohman, Marcus Michelen, Dhruv Mubayi}
\date{Source: arXiv:2603.16454v1 (CC BY 4.0)}
\begin{document}
\maketitle

\noindent\emph{Source and licence.} The largest $K_r$-free set of vertices in a random graph. Version arXiv:2603.16454v1 (CC BY 4.0), distributed under the Creative Commons Attribution 4.0 International licence, as recorded in the arXiv licence metadata for this exact version and shown on the abstract page https://arxiv.org/abs/2603.16454v1. Full rights notice: licenses/arxiv-2603.16454-CC-BY-4.0.txt.\par\smallskip

\noindent\textit{Editorial scope.} Counted unit: the Lemma whose source label is \texttt{lem:large-covers} in the article (source0.tex lines 548--551, source label \texttt{lem:large-covers}), delivered together with the article's own complete original proof of it (source0.tex lines 552--583, one \texttt{proof} environment of 32 lines). No mathematics is written, repaired or re-derived: statement and proof are the article's own text line for line. Required context retained uncounted: fact:Y-k-Delta (lines 407--409). Every definition, display and label of those lines is retained unchanged. Omitted from this package and not redistributed: 1--406, 410--547, 584--1163. Nothing inside a retained excerpt is omitted or altered: every retained body is delivered exactly as the article's lines are written, so no omission receipt of any kind accompanies this package. Citations: the retained excerpts carry no citation command at all, so no bibliography entry of the article is transcribed and no reference is redistributed. Reference adaptation: none was needed and none was made; every reference inside the delivered unit resolves inside it, so no unresolved reference or citation is produced. Printed slips kept literal: no printed word, symbol, hypothesis or number of the counted statement or of its proof was altered. Layout and class: the article's own class is \texttt{article} with packages including babel, inputenc, geometry, amsmath, amssymb, amsthm, graphicx, hyperref, enumitem, mathtools, cleveref, fullpage, tikz; the wrapper is the lane's pinned \texttt{amsart} editorial wrapper at 10pt on A4, which restates the theorem-like environments the retained text needs and adds the article's own macro definitions that the retained text needs (\texttt{\textbackslash E}, \texttt{\textbackslash EE}, \texttt{\textbackslash PP}, \texttt{\textbackslash cB}, \texttt{\textbackslash cC}, \texttt{\textbackslash cD}, \texttt{\textbackslash cZ}), with extra wrapper packages (bm, bbm, dsfont, xspace, cancel, mathtools, cleveref). The delivered numbering is the unit's own. Assets: no retained span contains a figure environment, a graphics-inclusion command, a PDF-page inclusion, a table or a TikZ picture; no image file is copied into this package, the package declares no asset dependency and no image file is redistributed with it. Licence: version arXiv:2603.16454v1 of this article is distributed under the Creative Commons Attribution 4.0 International licence, as recorded in the arXiv licence metadata for this exact version and shown on the abstract page https://arxiv.org/abs/2603.16454v1. A full rights notice is delivered at licenses/arxiv-2603.16454-CC-BY-4.0.txt. Characters: every retained excerpt is pure ASCII, so the delivered engine, XeLaTeX with its default Latin Modern text font, reports no missing character.
\begin{fact}\label{fact:Y-k-Delta}
    If $k = 2\log_2 n + O(\log \log n)$ then $$\binom{n}{k}\sum_{j = 1}^{k-1} \binom{k}{j}\binom{n-k}{k-j}2^{-2\binom{k}{2}+ \binom{j}{2}} \leq n^{-1 + o(1)}( \E [Y_k]^2 + \E [Y_k]) \,. $$
\end{fact}

\begin{lemma}\label{lem:large-covers}
    Suppose that $\E [Y_k] \in [n^{5/6},n^{6/5}]$.  Then for $i = O(\log\log n)$ we have 
    $$\PP(\cB_2 ) \leq  n^{-1/100}\quad \text{ and }\quad  \PP(\cB_3) \leq n^{-1/100}\,.$$
\end{lemma}

\begin{proof}
    Let $S \in \cZ_{k+1,i} \cup \cZ_{k+1, i+1}$ and $e \in \cD(S)$.  We begin with a union bound:  \begin{align} \label{unionb3}
        \PP(\cB_2) \leq (i \E[Z_{k+1,i}] + (i+1)\E[ Z_{k+1, i+1}]) \PP_{n-(k+1)}\left( |\cC(e)| \leq (3/4)^k \E[Y'_k] / 2 \right)
        %(\forall~T \in \cY_k, K_3 \subset G[e,T])
    \end{align}
   where we write $\PP_{n-(k+1)}$ for the distribution of the random graph on $ \binom{V}{2} \setminus \binom{S}{2}, e \in S$ and
   $ Y_k'$ is the number of independent sets of size $k$ that do not intersect $S$.
   
    
    To handle this probability, we will show by the second moment method that the number of independent sets that cover the disjoint edge $e$ is near its mean.  Define $N = n- (k +1)$.  Set $X_S = | \{ T \in \cC(e) : S \cap T = \emptyset\}|$ and note \begin{equation} \label{eq:E_N[X]}
        \E[ X_S] = (3/4)^k \binom{N}{k} 2^{-\binom{k}{2}} = (1 + o(1)) (3/4)^k \E [Y_k']\,.
    \end{equation}  We also can bound \begin{align*}
        \mathrm{Var}[X_S] - \E[X_S] &\le \binom{N}{k} \sum_{j = 1}^{k-1} \binom{k}{j} \binom{N - k}{k - j}2^{-2\binom{k}{2} + \binom{j}{2}} (3/4)^{2k-j} \\
        &= N^{-1 + o(1)} \left(\E[X_S]^2 + \E[X_S]\right)
    \end{align*}
    where the second line is by adapting \cref{fact:Y-k-Delta} and we recall $\E[X_S] = (1 + o(1))(3/4)^k \E[Y_k']$ by \eqref{eq:E_N[X]}.  By Chebyshev's inequality, we then see \begin{align*}
        \PP(X_S \leq (3/4)^k \E [Y_k'] / 2) \leq (4 + o(1))\frac{\mathrm{Var}[X_S]}{\E [X_S]^2} \leq  4 \left(N^{-1 + o(1)} + \frac{1 + o(1)}{\E[X_S]} \right) \leq n^{-1/99}
    \end{align*}
    where in the last line we used $5/6 + 2\log_2(3/4) \geq 1/99$. Consequently, 
   $$ \PP_{n-(k+1)}\left( |\cC(e)| \leq (3/4)^k \E[Y'_k] / 2 \right)=\PP(X_S \leq (3/4)^k \E [Y_k'] / 2) \leq n^{-1/99}.$$
   Finally, since $\EE[Z_{k+1, i}], \E[ Z_{k+1,i+1}]$, and $i$ are each $n^{o(1)}$, we obtain 
$\PP(\cB_2 ) \leq  n^{-1/100}$ from (\ref{unionb3}).
   
    For the bound on $\cB_3$ we may again union bound to see \begin{align*}
        \PP(\cB_3) &\leq (i \E[Z_{k+1,i}] + (i+1) \E[ Z_{k+1, i+1}]) \binom{N}{k} \sum_{j = 1}^{k-1} \binom{k}{j} \binom{N - k}{k - j}2^{-2\binom{k}{2} + \binom{j}{2}} (3/4)^{2k-j} \\
        &= n^{-1 + o(1)}( i\E[Z_{k+1,i}] + (i+1) \E[ Z_{k+1, i+1}]) \left(\E[X_S]^2 + \E[X_S]\right)\,. 
    \end{align*}
    As $ \E[ Y_k'] = (1+o(1)) \E[Y_k]$, recalling (\ref{eq:E_N[X]}) we have
    $$\E[X_S] = n^{2\log_2(3/4) + o(1)} \E[Y_k] \le n^{2\log_2(3/4) + 6/5+o(1)}
    \le n^{0.38 +o(1)}.
    $$ As $\E[Z_{k+1,i}], \E[Z_{k+1,i+1}]$ are each $n^{-1 + o(1)} \E[Y_k]$ this completes the proof.
\end{proof}

\end{document}
